\section{Limitations, Future Work and Conclusion}\label{sec:limits}
\textbf{Limitations.}
(i)~\emph{Sparse coverage of the gate operating points.} The remaining gate error comes from OFF transistors at rail-to-rail bias, where about 0.3\,\% of the training rows lie (Sec.~\ref{sec:errors:diag}); the project document fixes uniform sampling of $V_d,V_s$.
(ii)~\emph{The summation itself.} With perfect gate models the plain sum reaches only 99.16\,\% (Sec.~\ref{sec:circuit:ceiling}); about a quarter of the deficit comes from the input-pin currents and the rest, most plausibly, from gate inputs a few millivolts off the rails.
(iii)~\emph{Definitions.} Accuracy and gate leakage are defined by our assumptions (Sec.~\ref{sec:assumptions}).
(iv)~\emph{Single split, single seed.} All choices rest on one 80:10:10 split (seed 42) and one training seed; run-to-run standard deviations of the test median reach 0.40 for the ResNet, comparable to some of the differences used to choose between model types (e.g.\ the NMOS ON tie of 0.185, 0.186 and 0.186\,\%).
(v)~\emph{Data conventions.} PMOS OFF $i_g,i_b$ come from BSIM4 instance currents; six-digit CSV output limits what can be derived from node voltages; the corner stream depends on the ngspice build (Sec.~\ref{sec:data}).
(vi)~\emph{Scope.} Only DC operating points of fixed-size static gates and one adder in one technology are covered; no transient or loading effects.
% src: results/analysis/report_numbers/rail_point_row_counts.csv, results/partC/cla3_results.csv, results/analysis/loading_analysis/cla_ceiling_with_pins.csv, offrail_nodes.txt, archive/model_benchmark/seed_runs/summary_mean_std.csv

\textbf{Future work.} Within the specification the cheapest improvements are on the modelling side: more physics at the rails (a better analytical part of the hybrid model, features exposing the distance of the source to its rail), averaging several training seeds of the OFF ResNets, and a larger validation set to select by the near-rail error more reliably (it has only 90 near-rail rows per case). A learning curve (validation accuracy against training-set size) would show whether the OFF models are limited by data or by model form. Closing the summation gap needs the input-pin currents or predicted internal node voltages, i.e.\ a model beyond the four transistor models.

\textbf{Conclusion.} We built and checked a transistor $\rightarrow$ gate $\rightarrow$ circuit leakage estimator for the PTM \SI{22}{\nano\meter} HP technology. Four transistor models, one per case and chosen on a validation split from four model types, give test median errors of 0.03--0.59\,\% for the four terminal currents; assembled with Kirchhoff's current law on the supply node, they predict the leakage of five gates over 18 input cases and 200 PVT corners with a mean accuracy of 98.87\,\% (median 99.18\,\%) and, summed over 33 gates, the leakage of a 3-bit carry-lookahead adder over 64 vectors with a mean accuracy of 99.00\,\% (median 99.17\,\%). Against the 99\,\% target the gate level falls 0.13 points short and the circuit level meets it narrowly: the plain sum of isolated gate leakages caps the circuit accuracy at 99.16\,\% with perfect gate models, and the ML result benefits from a cancellation between the positive bias of the gate models and the negative bias of the summation. The remaining gate error is traced to OFF transistors at rail-to-rail bias, a corner that uniformly sampled training data cover with about 0.3\,\% of their rows.
% src: results/partA/chosen_models_test.csv (test medians 0.03--0.59), results/partB/gate_results.csv, results/partC/cla3_results.csv
